\section{Related work}
\label{sec:related}

\paragraph{Pairwise-comparison leaderboards.}
The Bradley--Terry model~\citep{bradley1952rank} and its MM/Newton estimators~\citep{hunter2004mm} underlie sports ratings and, more recently, LLM leaderboards such as Chatbot Arena, which fits BT scores to crowdsourced votes and reports bootstrap confidence intervals~\citep{chiang2024chatbot}.
MT-Bench~\citep{zheng2023judging} released expert and GPT-4 pairwise judgments of six chatbots and showed that GPT-4 agrees with humans about as often as humans agree with each other.
\citet{singh2025leaderboard} document structural issues of Arena (private testing, unequal sampling and data access across providers); our analysis is complementary, quantifying how much the published order depends on small parts of the data, and our finding that low-vote models are the cheapest to move gives a mechanism by which unequal sampling translates into unequal fragility.

\paragraph{Influence and dropping-data robustness.}
Influence functions~\citep{hampel1974influence,cook1977detection} approximate the effect of removing a data point; \citet{koh2017understanding} popularised them in machine learning.
\citet{broderick2020automatic} proposed the Approximate Maximum Influence Perturbation (AMIP): the smallest fraction of data whose removal can overturn a conclusion, approximated to first order, and showed that this sensitivity is governed by signal-to-noise and is not captured by standard errors.
\citet{huang2025dropping} applied this worst-case dropping idea to BT rankings of Chatbot Arena and MT-Bench, finding that dropping as little as $0.02\%$ of evaluations can change the top model, that MT-Bench is markedly more robust (which they attribute to expert annotators and curated prompts), and that crowdsourced and LLM-judge rankings are about equally sensitive.
We build on this line but differ in four ways: (i) we study three perturbation families (deletion, label flips, injection, including injection without control of match-ups) with every budget verified by exact refits rather than reported from the linear approximation; (ii) we derive and test a quantitative law, $k^\star\approx C\,z\sqrt{n_h}$, that collapses all pairs from all datasets onto one curve and \emph{explains} cross-design differences by signal-to-noise and per-entity sample size; (iii) we use controlled comparisons (identical battles labelled by crowd vs.\ six LLM judges; identical MT-Bench items labelled by experts vs.\ GPT-4; subsampling) to separate the effect of the labeller from that of the sampling design, which lets us test the expertise explanation directly; (iv) we include a non-LLM design (36 NBA seasons) and contrast worst-case with random perturbations and the bootstrap.

\paragraph{Manipulation of crowdsourced evaluation.}
\citet{min2025improving} show that a few hundred rigged votes can improve a model's Arena rank, using ``omnipresent'' strategies that exploit votes on battles not involving the target, because a target model appears in only about 1\% of new battles.
Our injection budgets give the idealised (free match-up) and the target-only (random match-up) end points of this threat model under BT scoring, and relate both to the leaderboard's own statistical precision.
